\section{Agentic Workflows}
\label{sec:agentic-workflows}

Agentic workflows are declarative specifications of multi-step processes that
agents execute autonomously within the governance constraints of the \kdd~model.

\subsection{Workflow Structure}

Each workflow definition contains:

\begin{itemize}
  \item \textbf{trigger}: the event or command that initiates the workflow,
  \item \textbf{steps}: an ordered sequence of agent actions,
  \item \textbf{inputs}: declared parameters and their sources,
  \item \textbf{outputs}: expected artefacts with metadata requirements,
  \item \textbf{governance\_checks}: validation rules applied before committing
    each output.
\end{itemize}

\subsection{Core Workflows}

The system defines five primary workflows:

\begin{description}
  \item[\workflow{pre-grand-prix}] Aggregates telemetry from previous sessions,
    retrieves historical setup patterns, and produces a race strategy brief.
  \item[\workflow{run-setup-what-if}] Simulates alternative setup configurations
    using the digital twin and ranks them by predicted lap time delta.
  \item[\workflow{generate-paper-section}] Produces a draft paper section from
    experiment results, metadata, and bibliography entries.
  \item[\workflow{create-new-feature}] Scaffolds a new feature across the
    relevant repositories following governance constraints.
  \item[\workflow{copilot-evidence-report}] Assembles a structured evidence
    report from active telemetry and knowledge base queries.
\end{description}

\subsection{Skill Reuse}

Workflows invoke reusable capabilities registered in \repo{04-skills-autoskills-registry}.
A skill encapsulates a validated, versioned agent capability. Invoking a registered
skill rather than a free prompt increases output consistency and audit coverage.
